\documentclass[11pt]{article}
\usepackage[T1]{fontenc}
\usepackage[utf8]{inputenc}
\usepackage{lmodern}
\usepackage[margin=1in]{geometry}
\usepackage{amsmath,amssymb,amsthm,mathtools}
\usepackage{microtype,booktabs,enumitem}
\usepackage{xcolor}
\usepackage[colorlinks=true,linkcolor=blue!45!black,citecolor=blue!45!black,urlcolor=blue!45!black]{hyperref}
\usepackage[nameinlink,noabbrev]{cleveref}
\hypersetup{pdftitle={Width-independent generalization in the strong-prior limit of Bayesian neural networks},pdfauthor={Dmitry Vaintrob; Kaarel Hanni; ChatGPT},pdfsubject={Strong-prior Bayesian generalization and norm-based compression}}
\newtheorem{theorem}{Theorem}[section]
\newtheorem{proposition}[theorem]{Proposition}
\newtheorem{lemma}[theorem]{Lemma}
\newtheorem{corollary}[theorem]{Corollary}
\theoremstyle{definition}
\newtheorem{definition}[theorem]{Definition}
\newtheorem{remark}[theorem]{Remark}
\newcommand{\R}{\mathbb R}
\newcommand{\E}{\mathbb E}
\newcommand{\PP}{\mathbb P}
\newcommand{\cH}{\mathcal H}
\newcommand{\cX}{\mathcal X}
\newcommand{\norm}[1]{\left\lVert#1\right\rVert}
\newcommand{\abs}[1]{\left|#1\right|}
\DeclareMathOperator{\err}{err}
\DeclareMathOperator{\Lip}{Lip}
\DeclareMathOperator{\sgn}{sgn}
\DeclareMathOperator{\diag}{diag}
\setlist{itemsep=2pt,topsep=5pt}
\setlength{\emergencystretch}{2em}
\title{Width-independent generalization in the strong-prior limit\ of Bayesian neural networks}
\author{Dmitry Vaintrob \and Kaarel H\"anni \and ChatGPT\thanks{AI assistant developed by OpenAI; contributed to drafting and mathematical checking.}}
\date{October 15, 2026}
\begin{document}
\maketitle
\begin{abstract}
We study noiseless Bayesian learning of Boolean functions by fully connected neural networks with bounded Lipschitz activations and a linear readout. The Gaussian prior uses layer-dependent variances chosen so that duplicating hidden neurons preserves a weighted squared norm. At every fixed depth, we prove a sample bound polynomial in the input dimension and the size of a bounded-weight network implementing the target, uniformly over all larger learner widths, in the limit of small prior variance at each fixed width. The proof has four ingredients: a posterior bound based on a finite set of relevant functions, a small-variance comparison of their Gaussian probabilities, a cloning construction, and a width-independent compression argument. Compression proceeds from the last hidden layer to the first, after pruning large edges row by row. This gives a polynomial equivalence between minimum weighted norm and bounded-weight network complexity. A separate tanh construction shows that the norm sublevel sets can grow exponentially with depth. The result concerns the strong-prior limit; it gives no width-uniform bound on how small the prior variance must be.
\end{abstract}

\section{Introduction}
An overparameterized neural network can fit a training sample in many ways. A generalization argument must therefore explain which of those fits the learning rule favors. For Bayesian learning, this question is about the prior distribution on functions induced by the parameter distribution. Even for a simple Gaussian parameter prior, the probability of an individual function can be difficult to estimate and can depend strongly on the network width.

We consider a regime in which a more tractable quantity determines the relevant functions: the smallest weighted squared parameter norm of a network implementing each function with a fixed output margin. When the Gaussian variance tends to zero at fixed width, a function with larger minimum norm has vanishing prior probability relative to a function with smaller minimum norm. Learning can then be controlled by counting the functions whose minimum norm is no larger than the target's. An absolute lower bound on the target's prior probability is unnecessary.

Our main result is that this count is at most exponential in a polynomial of the target's bounded-weight network complexity, with a polynomial independent of the learner's width. Two constructions establish the connection between these complexity measures. Cloning embeds a small target network into any larger width at bounded cost. Conversely, a network of bounded weighted norm can be compressed to polynomial width while preserving its Boolean function. The latter construction is the main technical step: it first removes large entries using thresholds proportional to their row norms, then samples hidden neurons in reverse layer order. Reverse order ensures that each concentration bound ranges over the already compressed next layer, rather than over the original width.

The proof develops the main argument of Vaintrob and H\"anni's lecture notes~\cite{notes}. We restrict attention to the small-variance limit at fixed width. We do not analyze the mean-field or Gaussian-process limits, noisy observations, gradient-based training, or generalization to a different input distribution. The depth is fixed in the polynomial learning guarantee. An appendix gives a rigorous tanh example explaining why depth dependence in the norm-counting argument cannot simply be discarded.

\paragraph{Relation to earlier work.}
The statistical step is a realizable Occam argument~\cite{blumer}, adapted to a posterior with a small tail outside a finite set. Norm-based capacity control independent of parameter count has a long history, including Bartlett~\cite{bartlett}, Neyshabur, Tomioka, and Srebro~\cite{neyshabur}, and Golowich, Rakhlin, and Shamir~\cite{golowich}. Their activation and norm assumptions differ, so we give a self-contained compression and counting proof. Buzaglo et al.~\cite{buzaglo} also analyze generalization of overparameterized networks through probabilities of fitting a teacher under random parameter sampling. Here the organizing quantity is a relative small-variance rate, and the theorem's width independence applies to sample complexity rather than to the required variance threshold.

\section{Setting and main result}\label{sec:setting}
Fix an input dimension $d\geq1$ and depth $L\geq2$. Write $\cX=\{-1,1\}^d$ and $\bar x=(1,x_1,\ldots,x_d)$. All $L-1$ hidden layers have width $N$. For a fixed activation $\phi:\R\to\R$ satisfying
\begin{equation}\label{eq:activation}
 |\phi(t)|\leq1,\qquad |\phi(t)-\phi(s)|\leq |t-s|,
\end{equation}
the network is
\[
 h^0(x)=\bar x,\qquad h^\ell(x)=\phi(W^\ell h^{\ell-1}(x))\quad(1\leq\ell<L),
 \qquad z_w(x)=W^Lh^{L-1}(x).
\]
The activation acts coordinatewise, and the readout $W^L$ is linear. There are no separate biases beyond the constant input coordinate; constant hidden neurons may carry an offset to later layers.

For a Boolean function $f:\cX\to\{-1,1\}$, let
\[
 U_{N,f}=\{w:f(x)z_w(x)>1/2\text{ for every }x\in\cX\}.
\]
These sets are open, disjoint, and indexed by a finite set of functions. Let $\cH_N=\{f:U_{N,f}\ne\varnothing\}$ and $U_N=\bigcup_{f\in\cH_N}U_{N,f}$. We only use widths for which $U_N$ is nonempty.

\begin{definition}[Norm and target complexity]
The weighted squared norm of a width-$N$ network is
\begin{equation}\label{eq:norm}
 \kappa_N(w)=\frac dN\norm{W^1}_F^2+
 \sum_{\ell=2}^{L-1}\norm{W^\ell}_F^2+N\norm{W^L}_2^2,
 \qquad
 \kappa_N(f)=\inf_{w\in U_{N,f}}\kappa_N(w).
\end{equation}
An empty sum is zero and an infimum over the empty set is $+\infty$. The bounded-weight network complexity $C_L(f)$ is the smallest integer $C\geq1$ such that $U_{C,f}$ contains a network all of whose weights have absolute value at most $C$.
\end{definition}

We assume $C_L(f)<\infty$ for the target. The learning theorem itself needs no non-polynomiality assumption beyond~\eqref{eq:activation}. To connect $C_L$ to ordinary Boolean circuits, we additionally take $\phi$ to be nonconstant; boundedness then makes it non-polynomial. Tanh is a running example.

Under the unconditioned Gaussian law $\mu_{N,\sigma}$, all weights are independent, with standard deviations
\begin{equation}\label{eq:prior}
 \sigma_1=\sqrt{N/d}\,\sigma,\qquad
 \sigma_\ell=\sigma\ (2\leq\ell<L),\qquad
 \sigma_L=\sigma/\sqrt N.
\end{equation}
Its density is proportional to $\exp[-\kappa_N(w)/(2\sigma^2)]$. Before observing data, condition this law on $U_N$. The induced function prior is
\[
 \pi_{N,\sigma}(g)=\frac{\mu_{N,\sigma}(U_{N,g})}{\mu_{N,\sigma}(U_N)}.
\]
This global margin conditioning is part of the model. It cannot be replaced silently by conditioning only on the observed training points.

Let $D$ be any distribution on $\cX$. Given $S_n=((x_i,f(x_i)))_{i=1}^n$ with independent $x_i\sim D$, the posterior $\Pi_{N,\sigma}(\cdot\mid S_n)$ is $\pi_{N,\sigma}$ restricted to functions consistent with all labels. Define
\[
 \err_D(g)=\PP_{x\sim D}\{g(x)\ne f(x)\},\qquad
 R_D(\Pi)=\E_{g\sim\Pi}\err_D(g).
\]
Thus $1-R_D(\Pi)$ is the average posterior probability of the correct label. A posterior-majority classifier, with any tie rule, has error at most $2R_D(\Pi)$.

\begin{theorem}[Strong-prior generalization]\label{thm:learning}
For every fixed $L\geq2$ there is a polynomial $p_L$ with the following property. Let $f$ have $C=C_L(f)<\infty$, and let $N\geq C$ be arbitrary. For every $\varepsilon,\delta\in(0,1)$ there is $\sigma_0>0$ such that, for all $0<\sigma<\sigma_0$, all distributions $D$ on $\cX$, and all
\begin{equation}\label{eq:sample}
 n\geq\frac{2}{\varepsilon}\left(p_L(d,C)+\log\frac1\delta\right),
\end{equation}
with probability at least $1-\delta$ over $S_n$,
\[
 R_D\bigl(\Pi_{N,\sigma}(\cdot\mid S_n)\bigr)\leq\varepsilon.
\]
The polynomial is independent of $N$. The threshold $\sigma_0$ may depend on $d,L,N,\phi,f,\varepsilon$, but not on $D,n$, or $\delta$.
\end{theorem}

Taking fixed $\varepsilon<10^{-3}$ and $\delta=e^{-C}$ gives the lecture notes' polynomial sample bound and average correct-label probability greater than $99.9\%$. The order of quantifiers matters: we first fix a width, however large, and then take the variance sufficiently small. The theorem provides neither a rate for $\sigma_0$ nor a simultaneous $N\to\infty$, $\sigma\to0$ theorem.

\section{Relevant functions and the small-variance limit}
\begin{lemma}[A finite relevant set]\label{lem:relevance}
Let $\pi$ be a prior on a countable class containing a target $f$ with $\pi(f)>0$. Suppose a finite set $A$ contains $f$, $|A|\leq M$, and
\[
 \pi(\cH\setminus A)\leq\tau\pi(f).
\]
For any $a,\delta\in(0,1)$, if $n\geq a^{-1}(\log M+\log(1/\delta))$, then the consistency posterior satisfies $R_D(\Pi)\leq a+\tau$ with probability at least $1-\delta$.
\end{lemma}
\begin{proof}
A function of error greater than $a$ survives $n$ independent examples with probability at most $e^{-an}$. A union bound over $A$ shows that, except on an event of probability at most $Me^{-an}$, every consistent member of $A$ has error at most $a$. On every training sample, the posterior denominator is at least $\pi(f)$, so the posterior mass outside $A$ is at most $\tau$. Its contribution to the posterior risk is at most $\tau$; the contribution inside $A$ is at most $a$.
\end{proof}

This lemma does not require $A$ to consist of the highest-prior functions. Nor does it require a unique simplest consistent function. Both points are useful when several functions have the same minimum norm.

\begin{lemma}[Gaussian rate for a function]\label{lem:rate}
At fixed $d,L,N$, for every $g\in\cH_N$,
\begin{equation}\label{eq:rate}
 \lim_{\sigma\downarrow0}2\sigma^2\log\mu_{N,\sigma}(U_{N,g})=-\kappa_N(g).
\end{equation}
Consequently, if $\kappa_N(g)>\kappa_N(f)$, then $\pi_{N,\sigma}(g)/\pi_{N,\sigma}(f)\to0$.
\end{lemma}
\begin{proof}
Use coordinates $u^\ell=\sqrt{c_\ell}W^\ell$, where $c_1=d/N$, $c_L=N$, and the other $c_\ell$ equal one. Then $\kappa_N(w)=\norm{u}_2^2$ and $u\sim\mathcal N(0,\sigma^2I_P)$, where $P$ is the parameter count. Let $V$ be the image of $U_{N,g}$ and $b=\inf_{u\in V}\norm{u}^2$.

For $0<\theta<1$, the Gaussian exponential-moment bound gives
\[
 \PP\{u\in V\}\leq\PP\{\norm{u}^2\geq b\}
 \leq\theta^{-P/2}\exp\left(-\frac{(1-\theta)b}{2\sigma^2}\right).
\]
Take the upper limit and then $\theta\downarrow0$. Conversely, for any $r>0$, openness and the definition of $b$ give a positive-volume ball $B\subset V$ on which $\norm{u}^2<b+r$. Integrating the density over $B$ gives
\[
 \PP\{u\in V\}\geq\operatorname{vol}(B)(2\pi\sigma^2)^{-P/2}
 \exp\left(-\frac{b+r}{2\sigma^2}\right).
\]
Take the lower limit and then $r\downarrow0$. This proves~\eqref{eq:rate}. Ratios of the conditioned function probabilities equal ratios of the unconditioned event probabilities, so their logarithmic rate is $-\kappa_N(g)+\kappa_N(f)$.
\end{proof}

For $A_{N,f}=\{g\in\cH_N:\kappa_N(g)\leq\kappa_N(f)\}$, finiteness of $\cH_N$ now implies
\begin{equation}\label{eq:tail}
 \frac{\pi_{N,\sigma}(\cH_N\setminus A_{N,f})}{\pi_{N,\sigma}(f)}\longrightarrow0.
\end{equation}
Equal-cost functions all remain in $A_{N,f}$; no claim about their limiting probability ratios is needed. Since there are finitely many targets at fixed $d,N$, a single sufficiently small variance can also enforce any prescribed tail ratio for every representable target at that width. This observation supplies no quantitative or width-uniform threshold.

\section{Cloning: a small target stays cheap}\label{sec:cloning}
\begin{lemma}[Cloning]\label{lem:cloning}
If a width-$C$ network implements $f$, then for every $N\geq C$ a width-$N$ network implements the same real-valued output and has weighted norm at most $2\kappa_C(w)$. If $N$ is a multiple of $C$, the norm is preserved exactly. In particular, for $C=C_L(f)$,
\begin{equation}\label{eq:cheap}
 \kappa_N(f)\leq B_L(d,C):=d(d+1)C^2+LC^4.
\end{equation}
\end{lemma}
\begin{proof}
Set $k=\lfloor N/C\rfloor$ and replace each hidden neuron by $k$ identical copies. Copy its first-layer incoming weights unchanged. For each later layer, divide every weight leaving a copied neuron by $k$; connect every destination copy in the same way. The copies compute identical activations, and their outgoing contributions sum to the original contribution. Pad unused positions with zero incident weights. This works even when $\phi(0)\ne0$, since the padded neurons have no outgoing effect.

The first-layer squared norm is multiplied by $k$, each internal squared norm is unchanged, and the readout squared norm is divided by $k$. After including the width-dependent coefficients, the first-layer contribution is multiplied by $kC/N\leq1$, and the readout contribution by $N/(kC)<2$. These factors both equal one if $N=kC$.

For a complexity-$C$ witness, the original first-layer contribution is at most $d(d+1)C^2$, the internal contribution at most $(L-2)C^4$, and the readout contribution at most $C^4$. Applying the separate factors above gives~\eqref{eq:cheap}.
\end{proof}

\section{Compression by reverse-order subsampling}\label{sec:compression}
\begin{theorem}[Compression]\label{thm:compression}
Let $d\geq1$, $L\geq2$, and let $\phi:\mathbb R\to[-1,1]$ be
$1$-Lipschitz. Suppose a width-$N$ network $W$ satisfies
$\kappa_N(W)\leq K$, where $K\geq1$. Set
\begin{equation}\label{eq:compression-parameters}
 \begin{gathered}
 \eta=\frac{1}{12(L-1)},\qquad
 \rho=\eta 2^{-L}K^{-(L-1)/2},\\
 m_L(d,K):=M=\max\left\{\lceil9K\rceil,
 \left\lceil256\rho^{-4}\bigl(d+\log(L/\rho)\bigr)\right\rceil\right\}.
 \end{gathered}
\end{equation}
There is a width-$M$ network $\widetilde W$ such that
\begin{equation}\label{eq:compression-conclusion}
 \max_{x\in\{-1,1\}^d}
 |z_{\widetilde W}(x)-z_W(x)|\leq\frac1{12},
 \qquad \kappa_M(\widetilde W)\leq4K.
\end{equation}
If $W$ implements a Boolean function $g$ with strict margin $1/2$,
then multiplying the readout of $\widetilde W$ by $3/2$ gives a
width-$M$ implementation of $g$ with strict margin $1/2$,
weighted norm at most $9K$, and every weight of absolute value at most
$M$. In particular, $C_L(g)\leq M$. For fixed $L$, $M$ is bounded by a
polynomial in $d$ and $K$.
\end{theorem}

\begin{proof}
We first truncate large entries and then sample the hidden neurons.
All norms of rows below are Euclidean, and all matrix norms with
subscript $F$ are Frobenius norms.

\emph{Truncation.}
For $2\leq s\leq L$, let $r_i^s=\|W^s_{i,*}\|_2$. In each nonzero
row, delete entries larger in absolute value than
\[
 \tau_i^s=\frac{2r_i^s}{\rho\sqrt N}.
\]
Leave zero rows unchanged, write $\overline W^s$ for the resulting
matrices, and set $\overline W^1=W^1$. Here and below $W^L$ is the
single readout row. Because all hidden activations have absolute value
at most one, deleting a row's large entries changes its preactivation
by at most
\[
 \sum_{j:\,|W^s_{ij}|>\tau_i^s}|W^s_{ij}|
 \leq \frac{(r_i^s)^2}{\tau_i^s}
 =\frac\rho2\sqrt N\,r_i^s.
\]
The vector of changes at an internal layer therefore has norm at most
$(\rho/2)\sqrt N\|W^s\|_F$. Every original or truncated internal
matrix has operator norm at most $\sqrt K$, and the readout has norm
at most $\sqrt{K/N}$. The activation is a contraction in Euclidean
norm. Thus truncating layer $s<L$ changes the output by at most
\[
 \frac\rho2\sqrt{NK}\,K^{(L-s-1)/2}\sqrt{K/N}
 =\frac\rho2 K^{(L-s+1)/2}\leq\frac\eta2.
\]
Truncating the readout changes it by at most
$(\rho/2)\sqrt K\leq\eta/2$. These bounds also hold when other
layers have already been truncated. Telescoping over the $L-1$
truncations gives total output error at most $1/24$.

\emph{Independent samples and three events.}
For each hidden layer $s$, independently choose a list
$I_s(1),\ldots,I_s(M)$ of independent uniform indices in
$\{1,\ldots,N\}$. Repeated indices represent separate copies of a
neuron. Define the fully sampled network by
\begin{align*}
 \widetilde W^1_{a,*}&=W^1_{I_1(a),*},\\
 \widetilde W^s_{a,b}&=\frac NM
      \overline W^s_{I_s(a),I_{s-1}(b)}
      &&(2\leq s\leq L-1),\\
 \widetilde W^L_b&=\frac NM\overline W^L_{I_{L-1}(b)}.
\end{align*}
For the analysis expose these independent lists in reverse order,
$I_{L-1},\ldots,I_1$. At the stage sampling layer $m$, all earlier
layers still have their original, truncated weights. Let
$\overline h_j^m(x)$ denote their exactly computed activations. The
estimated sum into a retained row $i$ of layer $m+1$ is
\[
 \frac NM\sum_{b=1}^M
 \overline W^{m+1}_{i,I_m(b)}\overline h^m_{I_m(b)}(x).
\]
Its expectation is the corresponding full sum over $N$ neurons.
Conditional on the later lists, its summands are independent and
bounded in absolute value by $(N/M)\tau_i^{m+1}$. Hoeffding's inequality~\cite{hoeffding}
shows that its deviation exceeds
$(\rho/2)\sqrt N r_i^{m+1}$ with probability at most
\[
 2\exp\left(-\frac{\rho^4 M}{32}\right).
\]
A zero row is estimated exactly. The bound is uniform in the retained
row, so it remains valid after averaging over the later lists.

Let $A$ be the event that all these row estimates, at all stages and
all $2^d$ inputs, satisfy their deviation bounds. At the last hidden
layer there is one output row; at every other stage there are $M$
retained rows. Consequently
\begin{equation}\label{eq:compression-hoeffding}
 \Pr(A^c)\leq 2^{d+1}LM\exp(-\rho^4M/32)\leq\frac1{100}.
\end{equation}
For completeness, put $a=d+\log L+\log(1/\rho)$ and
$M_0=\lceil256\rho^{-4}a\rceil$. Then
$M_0\leq257\rho^{-4}a$, and $u\exp(-\rho^4u/32)$ decreases for
$u\geq M_0$. Substituting $M_0$ and using $\log a\leq a-1$ bounds
the logarithm of the right side before its final inequality by
\[
 -(7-\log2)d-6\log L-3\log(1/\rho)+\log514-1
 \leq-\log100.
\]
This also proves the bound when $M$ was enlarged to be at least $9K$.

Next define the total retained row mass using the \emph{original}
row norms:
\[
 R=N\sum_{s=2}^{L-1}\sum_{a=1}^M(r^s_{I_s(a)})^2
      +N\|W^L\|_2^2.
\]
Its expectation is
\[
 \mathbb E R
 =M\sum_{s=2}^{L-1}\|W^s\|_F^2+N\|W^L\|_2^2
 \leq MK.
\]
By Markov's inequality the event $B=\{R\leq8MK\}$ has probability
at least $7/8$. In particular, on $B$ the retained row mass at each
individual stage is at most $8MK$.

Finally, direct calculation gives
\begin{align*}
 \mathbb E\left[\frac dM\|\widetilde W^1\|_F^2\right]
     &=\frac dN\|W^1\|_F^2,\\
 \mathbb E\|\widetilde W^s\|_F^2
     &=\|\overline W^s\|_F^2 &&(2\leq s\leq L-1),\\
 \mathbb E\left[M\|\widetilde W^L\|_2^2\right]
     &=N\|\overline W^L\|_2^2.
\end{align*}
For an internal matrix, independence of the two adjacent layers'
index lists makes each index pair uniform on $\{1,\ldots,N\}^2$;
the factor $(N/M)^2$ is canceled by the $M^2/N^2$ sampling factor.
Independence among different matrix entries is unnecessary.
Thus $\mathbb E\kappa_M(\widetilde W)\leq K$, and the event
$C=\{\kappa_M(\widetilde W)\leq4K\}$ has probability at least $3/4$.
On $C$ every sampled internal matrix has operator norm at most
$2\sqrt K$, and the sampled readout has norm at most $2\sqrt{K/M}$.

The events need not be independent:
\[
 \Pr(A\cap B\cap C)\geq1-\frac1{100}-\frac18-\frac14>0.
\]
Fix a collection of lists in their intersection.

\emph{Reverse-order telescoping.}
Let $F_m$ be the network with hidden layers $m,\ldots,L-1$ sampled
and layers $1,\ldots,m-1$ still original and truncated. Incoming
weights to a sampled layer select its rows, and outgoing weights from
a sampled layer are multiplied by $N/M$. This defines the intermediate
networks with different widths in different layers. In particular,
$F_L$ is the truncated network and $F_1$ is the fully sampled network.

When comparing $F_{m+1}$ and $F_m$, layer $m$ computes the original
truncated neuron functions before selection. The only new discrepancy
is in the sums into the retained neurons of layer $m+1$.
For $m=L-1$, event $A$ bounds the output discrepancy directly by
$(\rho/2)\sqrt K\leq\eta/2$. For $m<L-1$, put $t=L-m$.
Events $A$ and $B$ bound the Euclidean discrepancy in layer $m+1$'s
preactivation by $(\rho/2)\sqrt{8MK}$. All matrices \emph{after}
layer $m+1$ already equal their final sampled matrices, so their
operator norms are controlled by $C$. Propagating the discrepancy
therefore changes the output by at most
\begin{align*}
 \frac\rho2\sqrt{8MK}\,(2\sqrt K)^{t-2}
       2\sqrt{K/M}
 &=\rho\,2^{t-1/2}K^{t/2}\\
 &\leq\eta\,2^{-3/2}\leq\frac\eta2,
 \qquad 2\leq t\leq L-1.
\end{align*}
There are $L-1$ sampling stages, contributing at most $1/24$ in total.
Adding the truncation error proves the approximation bound in
\eqref{eq:compression-conclusion}; event $C$ proves its norm bound.

If the original network has strict margin $1/2$, the sampled network
has strict margin $5/12$. Multiplying its linear readout by $3/2$
restores strict margin $1/2$ and multiplies the weighted norm by at
most $9/4$. At width $M$ the resulting entrywise bounds are
\[
 |W^1_{aj}|\leq\sqrt{9KM/d},\qquad
 |W^s_{ab}|\leq3\sqrt K\ (1<s<L),\qquad
 |W^L_b|\leq3\sqrt{K/M}.
\]
Each is at most $M$ because $d\geq1$ and $M\geq9K\geq9$.
The definition of $C_L$ now applies. Finally, substituting
\eqref{eq:compression-parameters} gives
$\rho^{-4}=12^4(L-1)^4 16^L K^{2(L-1)}$, which establishes the stated
polynomial dependence at fixed depth.
\end{proof}

\section{Counting functions at bounded width}
\label{sec:packing}

The compression theorem reduces the counting problem to networks of bounded
width. We now count functions with low weighted-norm implementations at such
a width.
The argument uses the margin to separate representatives in parameter space.
It does not require the posterior distribution on parameters to be supported
in a bounded set.

Write $u_\ell=\sqrt{c_\ell}\,W^\ell$, where
$c_1=d/N$, $c_\ell=1$ for $2\leq\ell\leq L-1$, and $c_L=N$.
Concatenating these coordinates gives $\kappa_N(w)=\|u\|_2^2$.
The dimension of the parameter space is
\[
 P_N=N(d+1)+(L-2)N^2+N.
\]

\begin{lemma}[Parameter sensitivity]
\label{lem:parameter-sensitivity}
Suppose $|\phi|\leq1$ and $\operatorname{Lip}(\phi)\leq1$.
For $d\geq1$, $L\geq2$, and $R>0$, define
\[
 \Lambda_L(R,d)^2
 =\frac{d+1}{d}R^{2(L-1)}+\sum_{j=0}^{L-2}R^{2j}.
\]
For all rescaled parameters $u,v$ with $\|u\|_2,\|v\|_2\leq R$,
and every $x\in\{-1,1\}^d$,
\[
 |z_u(x)-z_v(x)|\leq\Lambda_L(R,d)\|u-v\|_2.
\]
\end{lemma}

\begin{proof}
Telescope through networks in which one layer at a time is replaced.
Every layer in each intermediate network comes from either $u$ or $v$,
and therefore its rescaled Frobenius norm is at most $R$.
The original first-layer Frobenius norm is at most $R\sqrt{N/d}$,
each intermediate-layer norm is at most $R$, and the readout norm is
at most $R/\sqrt N$.

Changing the first rescaled layer by $\Delta u_1$ changes its
preactivation vector by at most
$\sqrt{N/d}\sqrt{d+1}\,\|\Delta u_1\|_F$.
The subsequent activations are $1$-Lipschitz, so its contribution at the
output is at most
\[
 \sqrt{\frac{d+1}{d}}R^{L-1}\|\Delta u_1\|_F.
\]
For $2\leq\ell\leq L-1$, the input activation vector has norm at most
$\sqrt N$. Changing the $\ell$th layer therefore contributes at most
\[
 \sqrt N\,\|\Delta u_\ell\|_F
 R^{L-1-\ell}\frac{R}{\sqrt N}
 =R^{L-\ell}\|\Delta u_\ell\|_F
\]
at the output. Changing the rescaled readout contributes at most
$\|\Delta u_L\|_2$. Summing these contributions and applying
Cauchy--Schwarz gives the claim. Notice that the factors of $\sqrt N$
cancel, so this sensitivity bound is independent of the width.
\end{proof}

\begin{proposition}[Packing bound]
\label{prop:packing}
Under the hypotheses of Lemma~\ref{lem:parameter-sensitivity},
\[
 \#\{g:\kappa_N(g)\leq R^2\}
 \leq \bigl(1+2R\Lambda_L(R,d)\bigr)^{P_N}.
\]
\end{proposition}

\begin{proof}
For every counted function $g$, take a sequence of strict-margin
implementations whose costs tend to $\kappa_N(g)$.
Compactness gives a convergent subsequence in rescaled parameter space,
with limit $u_g$ satisfying $\|u_g\|_2\leq R$.
Continuity of the network on the finite input domain gives
$g(x)z_{u_g}(x)\geq1/2$ for every $x$.

If $g\ne g'$, their labels differ somewhere, so
$|z_{u_g}(x)-z_{u_{g'}}(x)|\geq1$ for some input $x$.
Lemma~\ref{lem:parameter-sensitivity} therefore gives
$\|u_g-u_{g'}\|_2\geq1/\Lambda_L(R,d)$.
The open balls of radius $1/(2\Lambda_L(R,d))$ centered at the $u_g$
are disjoint and lie inside the ball of radius
$R+1/(2\Lambda_L(R,d))$.
Comparing their $P_N$-dimensional volumes proves the bound.
\end{proof}



\section{Counting low-norm functions and proving the theorem}\label{sec:counting}
The preceding ingredients give both a uniform function count and a comparison of complexity measures. We record explicit quantities to make the polynomial dependence unambiguous, without attempting to optimize it.

For $K\geq1$, use the compression width $m_L(d,K)$ from \cref{thm:compression}, and put $m=m_L(d,K)$ and $R=3\sqrt K$. Write
\begin{equation}\label{eq:Q}
 \begin{split}
 P_m&=m(d+1)+(L-2)m^2+m,\\
 \Lambda_L(R,d)^2&=\frac{d+1}{d}R^{2(L-1)}+\sum_{j=0}^{L-2}R^{2j},\\
 Q_L(d,K)&=P_m\log\bigl(1+2R\Lambda_L(R,d)\bigr).
 \end{split}
\end{equation}
For fixed $L$, these quantities are bounded by polynomials in $d,K$.

\begin{proposition}[Uniform count and complexity comparison]\label{prop:count}
For every $N\geq1$ and $K\geq1$,
\begin{equation}\label{eq:count}
 \#\{g\in\cH_N:\kappa_N(g)\leq K\}\leq\exp Q_L(d,K+1).
\end{equation}
Every such $g$ has $C_L(g)\leq m_L(d,K+1)$. Together with cloning, this implies that at fixed depth, $C_L(f)$ and $\kappa_N(f)$ are polynomially related, with dependence on $d$ and with bounds independent of $N$, whenever $N\geq C_L(f)$.
\end{proposition}
\begin{proof}
For each counted $g$, choose a strict-margin representative of norm less than $K+1$. This small slack avoids assuming that the infimum defining $\kappa_N(g)$ is attained. Compression gives a width-$m_L(d,K+1)$ representative of the same Boolean function with weighted norm at most $9(K+1)$. The packing bound gives~\eqref{eq:count}.

The compression width is at least $9(K+1)$. The rescaled representative's first-layer entries have magnitude at most $3\sqrt{(K+1)m/d}\leq m$, its internal entries at most $3\sqrt{K+1}\leq m$, and its readout entries at most $3\sqrt{(K+1)/m}\leq m$. It is therefore a witness for $C_L(g)\leq m$. This implication applies to every represented $g$, without assuming in advance that $N\geq C_L(g)$. Cloning supplies the other direction.
\end{proof}

\begin{proof}[Proof of \cref{thm:learning}]
Let $A=A_{N,f}$. By cloning, $\kappa_N(f)\leq B_L(d,C)$. By Proposition~\ref{prop:count},
\[
 \log|A|\leq Q_L\bigl(d,B_L(d,C)+1\bigr).
\]
Choose a polynomial $p_L$ majorizing the right-hand side. By~\eqref{eq:tail}, there is $\sigma_0>0$ such that the prior mass outside $A$ is at most $(\varepsilon/2)\pi_{N,\sigma}(f)$ whenever $0<\sigma<\sigma_0$. Apply Lemma~\ref{lem:relevance} with $a=\tau=\varepsilon/2$. The resulting risk bound and sample requirement are exactly those in the theorem.
\end{proof}

\section{Boolean circuits as neural-network witnesses}
\label{sec:circuit-simulation}

This section relates the witness complexity $C_L$ to the usual size of a
Boolean circuit. We fix a finite Boolean gate basis with fan-in bounded by a
constant $q$. Circuit depth counts gate layers, whereas network depth $L$
counts weight matrices, including the final linear readout.

\begin{proposition}[Simulation of a fixed Boolean gate basis]
\label{prop:circuit-simulation}
Let $\phi$ be continuous and nonpolynomial. Fix a finite Boolean gate basis
of fan-in at most $q$. There is a constant $a$, depending only on $\phi$
and the basis, with the following property. If a Boolean function $f$ has
a circuit with $C\geq1$ gates and depth $D$, then for every $L\geq D+1$
with $L\geq2$, it has a depth-$L$ network of width at most $aC$ whose
weights have absolute value at most $a$. The network has only the constant
input $x_0=1$, no separate bias parameters, and its linear output equals
$f(x)$ on every Boolean input. In particular, after increasing $a$ if necessary,
\[
 C_L(f)\leq aC.
\]
\end{proposition}

\begin{proof}
We use the finite interpolation theorem of Pinkus
\cite[Theorem~5.1]{pinkus}: if $\phi$ is continuous and nonpolynomial,
then arbitrary values on $r$ distinct points can be interpolated by a
linear combination of $r$ functions of the form
$v\mapsto\phi(a\cdot v-\theta)$.
Apply it to the Boolean input patterns of each gate in the fixed basis,
and also to the unary identity gate. Choose one representation for each:
\[
 G(v)=\sum_{s=1}^{r_G}\lambda_{G,s}
       \phi(a_{G,s}\cdot v-\theta_{G,s}),
 \qquad v\in\{-1,1\}^{q_G},\qquad r_G\leq2^{q_G}.
\]
There are finitely many chosen representations, so all their coefficients
are bounded by a constant depending only on $\phi$ and the gate basis.

We supply affine shifts using a constant unit in each hidden layer.
Choose $\beta$ such that $c=\phi(\beta)\ne0$, which exists because
$\phi$ is not the zero function. In the first layer this unit receives
weight $\beta$ from $x_0=1$. In each later hidden layer it receives
weight $\beta/c$ from the preceding constant unit, and again outputs $c$.
A gate preactivation obtains its shift $-\theta$ through weight
$-\theta/c$ from this constant unit (and through $x_0$ in the first layer).

Carry each Boolean signal as the indicated linear combination of a
constant number of hidden activations. A gate in the next layer can read
this signal by substituting that linear combination into its
preactivation. At most $q$ signals occur in one gate, so the resulting
weights remain bounded by a constant determined by the finite dictionary.
Allocate disjoint blocks of units to distinct signals; this prevents any
uncontrolled accumulation of coefficients through shared encodings.

To account for edges skipping layers, relay every needed earlier signal
using the identity-gate representation. A circuit with $C$ gates and
fan-in at most $q$ references at most $qC$ distinct input coordinates,
and creates at most $C$ gate outputs. Thus at most $(q+1)C$ Boolean
signals need to be carried at any layer. Each uses a constant number of
units, in addition to the one constant unit, giving width at most $aC$.
This realizes the circuit after $D$ hidden layers. If $L-1>D$, relay
the final Boolean value through the remaining hidden layers by the same
identity dictionary. The linear readout extracts its exact value
$f(x)\in\{-1,1\}$, which in particular satisfies the required margin.
All constants used in these constructions are independent of $d,C,D,L$.
\end{proof}

\begin{remark}
Finite interpolation also shows that every Boolean function has a
finite $C_L$ for $L\geq2$: first interpolate its values on the $2^d$
Boolean inputs, and then relay the resulting exact Boolean value with
the identity dictionary. No approximation by near-identity layers is
needed. The circuit proposition above is a forward simulation statement;
it does not assert an efficient uniform procedure for evaluating an
arbitrary nonpolynomial activation or for constructing the interpolants.
\end{remark}


\section{Scope and depth dependence}
The argument separates two questions. The geometric question is how many Boolean functions can have low weighted norm; cloning and compression answer it uniformly in width. The probabilistic question is how much prior mass lies above a target's norm level; here we answer it only asymptotically at each fixed width. Quantitative control of this second question would be needed for a theorem at an explicitly specified nonzero variance or along a joint width--variance scaling.

Our complexity comparison is between a margin-constrained norm and a bounded-weight network implementation. It does not assert a matching lower bound on the sample complexity of learning every high-complexity target. It also does not identify the prior probability of a function with an exponential of its circuit description length: the Gaussian rate has a factor $1/(2\sigma^2)$, and equal rates need not have equal prefactors.

Depth enters compression through the amplification of errors by downstream operator norms. Proposition~\ref{prop:depth-lower-bound} shows, for tanh, that low-norm sublevel sets can indeed contain doubly exponentially many functions as depth increases in an appropriate joint parameter regime. This is a limitation of a uniform norm-sublevel counting strategy. It is not a lower bound on the sample complexity of the Bayesian learner for a circuit-simple target: the competing functions would also have to receive sufficient posterior mass.

Finally, every statistical statement here uses fresh inputs from the same distribution as the training inputs. Neither an adversarial online mistake bound nor a claim about total mistakes over an unbounded iid stream is proved. Those require arguments beyond the finite-sample posterior bound established above.

\appendix
% Standalone appendix fragment; uses only standard amsmath/amssymb.
\section{A depth-dependent lower bound for the norm ball}
\label{app:depth-lower-bound}
The fixed-depth restriction in the counting argument has mathematical content.
For the activation $\phi=\tanh$, bounded weighted norm can encode many functions
when depth grows.  This is a statement about the norm ball, and does not by
itself give a lower bound for learning a circuit-simple target.

\begin{proposition}\label{prop:depth-lower-bound}
There is a universal constant $A$ with the following property.  Let $L\geq4$,
$1\leq q\leq d$, and $N\geq2^{q+1}$.  Every Boolean function depending on a
specified set of $q$ input coordinates has a width-$N$, depth-$L$ network with
$\tanh$ hidden activations, linear readout, output exactly in $\{-1,1\}$,
and
\[
 \kappa_N(w)\leq A\left[dq^2\log^2(2q)
             +(L-2)2^{2q/(L-2)}\right].
\]
Consequently, whenever the right-hand side is at most $K$, the norm ball
contains at least $2^{2^q}$ distinct Boolean functions.
\end{proposition}

\begin{proof}
We first work at width $N_0=2M$, where $M=2^q$, and write $m=L-3$.
For each $p\in\{-1,1\}^q$, introduce a first-layer unit
\[
 h_p(x)=\tanh\!\left(\gamma(\langle p,x\rangle-(q-1))\right),
 \qquad \gamma=\tfrac12\log(16q),
\]
where only the selected coordinates are used.  The other $M$ first-layer
units are constants $b=\tanh(1)$, obtained from the constant input coordinate.
The first-layer contribution to the weighted norm is exactly
\[
 \frac d2\left[\gamma^2\bigl(q+(q-1)^2\bigr)+1\right]
 =O\!\left(dq^2\log^2(2q)\right).
\]

Index the $q$-cube by its vertices and define the $M\times M$ matrix
$T_{xp}=(1+h_p(x))/2$.  Its diagonal differs from $1$ by at most
$e^{-2\gamma}$.  If the Hamming distance from $x$ to $p$ is $r\geq1$, then
$0<T_{xp}\leq e^{2\gamma-4\gamma r}$.  Hence
\[
 \|T-I\|_\infty
 \leq e^{-2\gamma}+e^{2\gamma}
            \left[(1+e^{-4\gamma})^q-1\right]<\tfrac14.
\]
Here $\|\cdot\|_\infty$ is the maximum absolute row sum; the last inequality
follows from $e^{-2\gamma}=(16q)^{-1}$ and
$(1+u)^q-1\leq qu e^{qu}$.  Thus $T$ is invertible and
$\|T^{-1}\|_\infty\leq4/3$.

Let $f$ denote the vector of desired labels, put $v=T^{-1}f$, and fix an
amplitude $0<t_0\leq s=3/4$.  One neuron in the second layer uses weights
$a_p=\frac12\operatorname{atanh}(t_0)v_p$ on the bump units and the common
weight $(\sum_p a_p)/(Mb)$ on the constant units.  Its preactivation is
\[
 \sum_p a_p h_p(x)+\sum_p a_p
   =\operatorname{atanh}(t_0)f(x),
\]
so its activation is exactly $t_0f(x)$.  Its squared weight norm is
$O(Mt_0^2)$, uniformly for $t_0\leq s$.  All other neurons in this layer
have zero weights and activation zero.

Choose
\[
 t_0=sM^{-(m-1)/(2(m+1))},\qquad
 R=(s/t_0)^{1/m},\qquad
 n_j=\left\lceil N_0^{j/m}\right\rceil,\quad 0\leq j\leq m.
\]
The $j$th of the next $m$ layers has $n_j$ active neurons, all computing
$t_jf$ with $t_j=t_0R^j$.  Given the preceding $n_{j-1}$ copies of
$t_{j-1}f$, assign each incoming weight
\[
 \frac{\operatorname{atanh}(t_j)}{n_{j-1}t_{j-1}}.
\]
This realizes the claimed activation exactly.  Since $0<t_j\leq s$,
$\operatorname{atanh}(t_j)\leq c t_j$ for a universal constant $c$.
Also $n_j/n_{j-1}\leq2N_0^{1/m}$.  The norm contribution of this layer is
therefore at most
\[
 \frac{n_j}{n_{j-1}}
 \left(\frac{\operatorname{atanh}(t_j)}{t_{j-1}}\right)^2
 \leq 2c^2N_0^{1/m}R^2
 =O\!\left(M^{2/(m+1)}\right).
\]
The second-layer cost is of this same order because
$Mt_0^2=s^2M^{2/(m+1)}$.  The last hidden layer has $N_0$ copies of $sf$;
readout weights $1/(N_0s)$ give output $f$ and weighted readout cost
$s^{-2}$.  Summing gives the asserted bound at width $N_0$.

Finally, for any $N\geq N_0$, replace every hidden neuron by either
$k=\lfloor N/N_0\rfloor$ or $k+1$ identical copies, choosing the
multiplicities in each layer to sum to $N$.  Divide each outgoing weight
by the multiplicity of its source neuron.  This preserves the function
exactly.  For an inner-layer edge, its contribution to squared norm is
multiplied by the ratio of its target and source multiplicities, at most
$(k+1)/k\leq2$.  The first- and last-layer weighted costs are likewise
multiplied by at most $2$, using their coefficients $d/N$ and $N$.
Absorbing this factor into $A$ proves the result.
\end{proof}

For example, taking $q=d\geq2$ and $L=d+2$ gives $2^{2^d}$ functions at weighted
norm $O(d^3\log^2(2d))$.  Thus a bound on the logarithm of the norm-ball
cardinality that holds uniformly in the width $N$ and is polynomial jointly
in $d$, $K$, and $L$ is impossible.
This conclusion concerns uniform counting; it does not establish that the
posterior needs exponentially many samples for a particular simple target.


\begin{thebibliography}{9}
\bibitem{notes}
Dmitry Vaintrob and Kaarel H\"anni.
\emph{A Generalization Result for Overparametrized Neural Net Bayesian Learning}.
ILIAD lecture slides, 2026.
\url{https://kaarelh.github.io/ai/nnbayes-generalization.pdf}.

\bibitem{blumer}
Anselm Blumer, Andrzej Ehrenfeucht, David Haussler, and Manfred K. Warmuth.
Occam's razor. \emph{Information Processing Letters}, 24(6):377--380, 1987.
\href{https://doi.org/10.1016/0020-0190(87)90114-1}{doi:10.1016/0020-0190(87)90114-1}.

\bibitem{bartlett}
Peter L. Bartlett. The sample complexity of pattern classification with neural networks: the size of the weights is more important than the size of the network.
\emph{IEEE Transactions on Information Theory}, 44(2):525--536, 1998.
\href{https://doi.org/10.1109/18.661502}{doi:10.1109/18.661502}.

\bibitem{neyshabur}
Behnam Neyshabur, Ryota Tomioka, and Nathan Srebro.
Norm-based capacity control in neural networks.
\emph{Proceedings of COLT}, PMLR 40:1376--1401, 2015.
\url{https://proceedings.mlr.press/v40/Neyshabur15.html}.

\bibitem{golowich}
Noah Golowich, Alexander Rakhlin, and Ohad Shamir.
Size-independent sample complexity of neural networks.
\emph{Information and Inference}, 9(2):473--504, 2020.
\href{https://doi.org/10.1093/imaiai/iaz007}{doi:10.1093/imaiai/iaz007}.

\bibitem{buzaglo}
Gon Buzaglo, Itamar Harel, Mor Shpigel Nacson, Alon Brutzkus, Nathan Srebro, and Daniel Soudry.
How uniform random weights induce non-uniform bias: typical interpolating neural networks generalize with narrow teachers.
\emph{Proceedings of ICML}, PMLR 235:5035--5081, 2024.
\url{https://proceedings.mlr.press/v235/buzaglo24a.html}.

\bibitem{pinkus}
Allan Pinkus. Approximation theory of the MLP model in neural networks.
\emph{Acta Numerica}, 8:143--195, 1999.
\href{https://doi.org/10.1017/S0962492900002919}{doi:10.1017/S0962492900002919}.

\bibitem{hoeffding}
Wassily Hoeffding. Probability inequalities for sums of bounded random variables.
\emph{Journal of the American Statistical Association}, 58(301):13--30, 1963.
\href{https://doi.org/10.1080/01621459.1963.10500830}{doi:10.1080/01621459.1963.10500830}.
\end{thebibliography}
\end{document}
